\documentclass[11pt]{article}
\usepackage[margin=1.1in]{geometry}
\usepackage{amsmath,amssymb,amsthm}
\usepackage[T1]{fontenc}
\usepackage[colorlinks=true,linkcolor=blue,citecolor=blue]{hyperref}

\newtheorem{theoremA}{Theorem}
\renewcommand{\thetheoremA}{\Alph{theoremA}}
\newtheorem{theorem}{Theorem}[section]
\newtheorem{lemma}[theorem]{Lemma}
\newtheorem{proposition}[theorem]{Proposition}
\newtheorem{corollary}[theorem]{Corollary}
\theoremstyle{definition}
\newtheorem{hypothesis}[theorem]{Hypothesis}
\newtheorem{example}[theorem]{Example}
\theoremstyle{remark}
\newtheorem{remark}[theorem]{Remark}

\newcommand{\R}{\mathbb{R}}
\newcommand{\N}{\mathbb{N}}
\newcommand{\Z}{\mathbb{Z}}
\newcommand{\Sone}{S^{1}}
\newcommand{\Lb}{L_{b}}
\newcommand{\Ab}{A_{b}}
\newcommand{\Tb}{T_{b}}
\newcommand{\Td}{T_{d}}
\newcommand{\expb}{\exp_{b}}
\newcommand{\frc}[1]{\{#1\}}
\newcommand{\Ht}{\widetilde H}
\newcommand{\Kt}{\widetilde K}
\DeclareMathOperator{\slog}{slog}
\DeclareMathOperator{\Var}{Var}

\title{The Mixed-Base Tetration Phase Map Is a Circle Diffeomorphism\\[2mm]
\large Paper II of the mixed-base tetration series: companion to ``A Phase Law for Mixed-Base Tetration'' \cite{J}}
\author{Janis Justus}
\date{July 2026}

\begin{document}
\maketitle

\begin{abstract}
Proposition~6.6 of \cite{J} derives the exact inverse mean antisymmetry
$\mu_{b,d}=-\mu_{d,b}$ for the mixed-base tetration phase under the
\emph{assumption} that $H(\theta)=\theta+\Phi_{b,d}(\theta)$ is an
orientation-preserving $C^{1}$ circle diffeomorphism. This paper removes
that assumption in four graded steps. \textbf{(A)}~Given the phase law of
\cite{J} for the two ordered pairs $(b,d)$ and $(d,b)$, the lift
$\Ht_{b,d}=\mathrm{id}+\Phi_{b,d}$ is unconditionally an
orientation-preserving circle \emph{homeomorphism} with inverse
$\Ht_{d,b}$, and the mixed-base phase maps form a groupoid in
$\mathrm{Homeo}_{+}(\Sone)$. \textbf{(B)}~Already at this topological
level, $\mu_{b,d}=-\mu_{d,b}$ holds exactly, by a Riemann--Stieltjes
substitution---no smoothness whatsoever. \textbf{(C)}~A minimal-regularity
principle: if \emph{both} phases are $C^{1}$ on the circle, the phase
maps are automatically mutually inverse $C^{1}$ diffeomorphisms; no
positivity estimate is needed, and an explicit example shows that
one-sided $C^{1}$ regularity would not suffice. \textbf{(D)}~Under an
explicit $C^{1}$ strengthening of the channel hypotheses of
\cite{J}---differentiable strict monotonicity of the branches and
quantitative positivity of the logarithmic channel, both satisfied by the
regular Kneser/Paulsen--Cowgill branch---the forward phase alone is
$C^{1}$ with $\Ht'>0$. The engine of (D) is an \emph{exact} product
identity for the derivative of the logarithmic tails,
\[
D_n'(\theta)\;=\;\alpha\,\Td'(\theta)\prod_{j=1}^{n-1}
\frac{\alpha\,\Td(j+\theta)}{\expb^{\circ j}\!\big(D_n(\theta)\big)},
\qquad \alpha=\log_b d,
\]
in which all tower-sized factors cancel; the summability of reciprocal
tower minima that drives the limit is itself automatic, since taking
infima in the tower recursion yields $s_{j+1}=d^{\,s_j}$. In the limit,
\[
\Ht_{b,d}'(\theta)\;=\;\frac{\alpha\,\Td'(\theta)}{\Tb'(\Ht_{b,d}(\theta))}
\prod_{j=1}^{\infty}\frac{\alpha\,\Td(j+\theta)}{\Tb\!\big(j+\Ht_{b,d}(\theta)\big)}\;>\;0 .
\]
All statements are validated numerically, including an end-to-end check
on explicit $C^{1}$ tetration branches.
\end{abstract}

\section{Introduction and statement of results}\label{sec:intro}

We keep the notation of \cite{J}. For admissible bases $a$, $T_a$ denotes
the regular tetration branch with $T_a(0)=1$ and $T_a(x+1)=a^{T_a(x)}$, and
$A_a=T_a^{-1}=\slog_a$ its Abel inverse. For an ordered pair $(b,d)$ one
sets $F_{b,d}=\Ab\circ\Td$, $\alpha=\log_b d$, $\Lb=\log_b$, and, for
$x=n+\theta$ with $n\in\N$, $\theta\in[0,1]$,
\[
D_n(\theta)=D_n^{b,d}(\theta):=\Lb^{\circ n}\big(\Td(n+\theta)\big).
\]
Under Hypotheses~2.2 and 2.4 of \cite{J}, Theorems~5.6 and 5.7 of \cite{J}
give the uniform limit $D=D_{b,d}=\lim_n D_n$ on $[0,1]$ and the phase law
\begin{equation}\label{eq:phaselaw}
\slog_b(\Td(x))-x \;=\; \Phi_{b,d}(\frc{x})+o(1),
\qquad x\to+\infty,
\qquad
\Phi_{b,d}(\theta)=\Ab(D(\theta))-\theta ,
\end{equation}
with $\Phi_{b,d}$ continuous and, by Lemma~5.8 of \cite{J}, well defined on
the circle. Write $\mu_{b,d}=\int_0^1\Phi_{b,d}$ and
$P_{b,d}=\Phi_{b,d}-\mu_{b,d}$.

Proposition~6.6 of \cite{J} proves the exact inverse antisymmetry
$\mu_{b,d}=-\mu_{d,b}$ under the standing \emph{assumption} that
\[
H(\theta)=\theta+\Phi_{b,d}(\theta)
\]
is an orientation-preserving $C^{1}$ circle diffeomorphism in the lifted
sense $H(\theta+1)=H(\theta)+1$, $H'>0$. The purpose of this paper is to
prove that assumption. Throughout, for a $1$-periodic continuous
$\Phi:\R\to\R$ we write
\[
\Ht(x):=x+\Phi(x),\qquad x\in\R,
\]
for the associated \emph{lift}; $\Ht(x+1)=\Ht(x)+1$ holds automatically.
Our results form a ladder of four statements, graded by the regularity
they consume: a topological theorem, a measure-level theorem, a formal
$C^{1}$ principle, and an analytic $C^{1}$ theorem.

\begin{theoremA}[Homeomorphism; no smoothness required]\label{thm:A}
Assume the phase-law package \textup{(P)} of Section~\ref{sec:setting}
holds for both ordered pairs $(b,d)$ and $(d,b)$ (in particular, this is
the case under Hypotheses~2.2 and 2.4 of \cite{J} for both pairs). Then
$\Ht_{b,d}=\mathrm{id}+\Phi_{b,d}$ is a strictly increasing homeomorphism
of $\R$ commuting with the unit translation, and
\[
\Ht_{b,d}^{-1}=\Ht_{d,b}.
\]
Consequently $\Ht_{b,d}$ descends to an orientation-preserving
homeomorphism $H_{b,d}$ of the circle $\Sone=\R/\Z$ with
$H_{b,d}^{-1}=H_{d,b}$. More generally, whenever the package \textup{(P)}
holds for the pairs $(p,q),(q,r),(p,r)$, the lifts compose exactly,
$\Ht_{p,r}=\Ht_{p,q}\circ\Ht_{q,r}$; the mixed-base phase maps form a
groupoid inside $\mathrm{Homeo}_{+}(\Sone)$.
\end{theoremA}

\begin{theoremA}[Exact inverse mean antisymmetry without smoothness]\label{thm:B}
Under the hypotheses of Theorem~\textup{\ref{thm:A}},
\[
\mu_{b,d}=-\mu_{d,b}
\]
holds exactly. In particular, the $C^{1}$ diffeomorphism assumption in
Proposition~6.6 of \cite{J} can be dropped entirely: the conclusion of
that proposition is valid whenever the phase law holds for both ordered
pairs.
\end{theoremA}

Theorem~\ref{thm:B} explains why the inverse-pair sums in Table~1 of
\cite{J} vanish to $10^{-9}$--$10^{-11}$: they are exactly zero in the
limit, and the observed residuals are numerical error.

\begin{theoremA}[Minimal-regularity principle]\label{thm:C}
Let $\Phi,\Psi:\R\to\R$ be $1$-periodic functions of class $C^{1}$
satisfying the inverse phase equations
\begin{equation}\label{eq:invpair}
\Psi\big(x+\Phi(x)\big)=-\Phi(x),
\qquad
\Phi\big(u+\Psi(u)\big)=-\Psi(u),
\qquad x,u\in\R .
\end{equation}
Then $\Ht=\mathrm{id}+\Phi$ and $\Kt=\mathrm{id}+\Psi$ satisfy
$\Kt\circ\Ht=\Ht\circ\Kt=\mathrm{id}$ and $\Ht'>0$, $\Kt'>0$ everywhere;
they are lifts of mutually inverse orientation-preserving $C^{1}$ circle
diffeomorphisms. Under \textup{(H1)}+\textup{(H2)} of
Section~\ref{sec:setting}, the equations \eqref{eq:invpair} hold
automatically for $\Phi=\Phi_{b,d}$, $\Psi=\Phi_{d,b}$
(Lemma~\ref{lem:comp}), so two-sided $C^{1}$ regularity of the phases
already implies the full diffeomorphism conclusion, with no separate
positivity assumption. One-sided $C^{1}$ regularity does not suffice
(Example~\ref{ex:counter}).
\end{theoremA}

\begin{theoremA}[Analytic $C^{1}$ theorem]\label{thm:D}
Assume, for the ordered pair $(b,d)$, Hypotheses~2.2 and 2.4 of \cite{J}
together with the strengthenings \textup{(S)} and \textup{(R$'$)} of
Section~\ref{sec:setting}. Then:
\begin{enumerate}
\item[\textup{(i)}] $D_n\to D$ in $C^{1}([0,1])$, $D\in C^{1}([0,1])$,
and $D'>0$ with the absolutely and uniformly convergent product
representation
\begin{equation}\label{eq:Dprime}
D'(\theta)\;=\;\alpha\,\Td'(\theta)\,
\prod_{j=1}^{\infty}\frac{\alpha\,\Td(j+\theta)}{E_j(\theta)},
\qquad
E_j(\theta):=\expb^{\circ j}\big(D(\theta)\big),
\end{equation}
where $\expb(t)=b^{t}$.
\item[\textup{(ii)}] $\Phi_{b,d}$ extends to a $1$-periodic $C^{1}$
function on $\R$, and the lift $\Ht_{b,d}=\mathrm{id}+\Phi_{b,d}$
satisfies $\Ht_{b,d}'=(A_b'\circ D)\,D'>0$; explicitly,
\begin{equation}\label{eq:Hprime}
\Ht_{b,d}'(\theta)\;=\;\frac{\alpha\,\Td'(\theta)}{\Tb'\big(\Ht_{b,d}(\theta)\big)}\,
\prod_{j=1}^{\infty}
\frac{\alpha\,\Td(j+\theta)}{\Tb\big(j+\Ht_{b,d}(\theta)\big)}\;>\;0 .
\end{equation}
\item[\textup{(iii)}] $H_{b,d}$ is an orientation-preserving $C^{1}$
circle diffeomorphism; in particular the standing assumption of
Proposition~6.6 of \cite{J} is a theorem under these hypotheses. If in
addition the package \textup{(P)} holds for $(d,b)$, then
$H_{d,b}=H_{b,d}^{-1}$ is automatically $C^{1}$ as well, i.e.\
$\Phi_{d,b}\in C^{1}(\Sone)$ comes for free.
\item[\textup{(iv)}] The residual converges in the $C^{1}$ topology:
with $R_n(\theta):=F_{b,d}(n+\theta)-(n+\theta)$ one has $R_n\to\Phi_{b,d}$
in $C^{1}([0,1])$; equivalently, $F_{b,d}'(n+\theta)\to
\Ht_{b,d}'(\theta)$ uniformly on $[0,1]$.
\end{enumerate}
\end{theoremA}

The engine of Theorem~\ref{thm:D} is the exact identity of
Lemma~\ref{lem:product} below, in which the mismatch constant $\alpha$
cancels against the first logarithmic peel and \emph{all} tower-sized
quantities cancel in closed form, leaving a product whose factors converge
to $1$ at the hyper-fast rate $O(1/s_j)$, $s_j=\inf_\theta\Td(j+\theta)$.
The summability of $1/s_j$ that drives the limit is not an extra
hypothesis: taking infima over the phase in the tower recursion gives
$s_{j+1}=d^{\,s_j}$ exactly (Lemma~\ref{lem:growth}), so the qualitative
divergence $s_j\to\infty$ of \cite[Hyp.~2.4]{J} self-improves to geometric
growth. As a by-product one obtains the sharp value-level comparison
\[
\Tb\big(j+\Ht_{b,d}(\theta)\big)=\alpha\,\Td(j+\theta)\big(1+O(1/s_j)\big)
\]
(Remark~\ref{rem:onelevel}): the base-mismatch factor $\alpha$ sits
exactly one tower level down.

The four statements interlock as follows. Theorem~\ref{thm:D}, applied to
the single pair $(b,d)$, already yields the diffeomorphism property
including positivity. Theorem~\ref{thm:C} shows that once \emph{both}
phases are known to be $C^{1}$---by Theorem~\ref{thm:D} applied to both
pairs, or by any future route such as strip analyticity
(Remark~\ref{rem:Cr})---positivity is automatic and free of channel
estimates. Example~\ref{ex:counter} delimits Theorem~\ref{thm:C}:
$C^{1}$ regularity of the forward phase alone is compatible with a
critical point, so the analytic positivity of Theorem~\ref{thm:D} is
genuinely needed in the one-sided situation. Theorems~\ref{thm:A} and
\ref{thm:B} need no smoothness at all.

\section{Setting and hypotheses}\label{sec:setting}

\subsection{Branch frame}
Throughout, $b,d>1$. We use the following normalization of the regular
branches, satisfied by the Kneser/Paulsen--Cowgill tetration for real
bases above $e^{1/e}$ \cite{Kneser,PC}.

\begin{itemize}
\item[\textup{(B1)}] $\Tb:(-2,\infty)\to\R$ is a continuous, strictly
increasing bijection onto $\R$ with $\Tb(0)=1$ and
$\Tb(x+1)=b^{\Tb(x)}$ for $x>-2$. Its inverse
$\Ab=\slog_b:\R\to(-2,\infty)$ then satisfies the \emph{global Abel
equation}
\begin{equation}\label{eq:abel}
\Ab(b^{y})=\Ab(y)+1\qquad\text{for all }y\in\R .
\end{equation}
\item[\textup{(B2)}] $\Td:[0,\infty)\to(0,\infty)$ is continuous with
$\Td(0)=1$ and $\Td(x+1)=d^{\Td(x)}$ for $x\ge 0$.
\end{itemize}

Equation \eqref{eq:abel} follows from (B1): if $x=\Ab(y)$ then
$b^{y}=b^{\Tb(x)}=\Tb(x+1)$, so $\Ab(b^{y})=x+1$. Note also
$\Tb(-1)=0$ and $\Tb>0$ exactly on $(-1,\infty)$, directly from the
functional equation.

\subsection{The phase-law package}
For an ordered pair $(p,q)$ of admissible bases we say that the
\emph{phase-law package} (P) holds if the conclusions of Theorems~5.6,
5.7 and Lemma~5.8 of \cite{J} are available for that pair, namely:
$D_n^{p,q}\to D_{p,q}$ uniformly on $[0,1]$; the residual satisfies
\eqref{eq:phaselaw} with the $o(1)$ uniform in the phase; and
$\Phi_{p,q}$ is continuous on $[0,1]$ with $\Phi_{p,q}(1)=\Phi_{p,q}(0)$,
so that its $1$-periodic extension (still denoted $\Phi_{p,q}$) is
continuous on $\R$. By \cite{J} this package holds under Hypotheses~2.2
and 2.4 of \cite{J} for the pair $(p,q)$; the endpoint identity is
re-proved in sharpened form in Lemma~\ref{lem:endpoint} below. For the
diagonal pairs the package is trivial: $F_{p,p}=\mathrm{id}$ exactly, so
$\Phi_{p,p}\equiv 0$.

We label:
\begin{itemize}
\item[\textup{(H1)}] the package (P) holds for $(b,d)$;
\item[\textup{(H2)}] the package (P) holds for $(d,b)$.
\end{itemize}

\subsection{$C^{1}$ strengthenings}
The following two conditions upgrade the channel hypotheses of \cite{J}
to the $C^{1}$ category. Recall $s_j:=\inf_{0\le\theta\le1}\Td(j+\theta)$
from Hypothesis~2.4 of \cite{J}.

\begin{itemize}
\item[\textup{(S)}] \emph{(Smooth strictly increasing branches.)}
$\Tb\in C^{1}((-2,\infty))$ with $\Tb'>0$, and
$\Td\in C^{1}([0,\infty))$ with $\Td'>0$. Consequently
$\Ab\in C^{1}(\R)$ with
\begin{equation}\label{eq:Abprime}
\Ab'(y)=\frac{1}{\Tb'(\Ab(y))}>0 ,\qquad y\in\R .
\end{equation}
\item[\textup{(R$'$)}] \emph{(Quantitative real channel.)} There is
$N_0\in\N$ such that for all $n\ge N_0$ and all $\theta\in[0,1]$ the
iterates
\[
W_k^{(n)}(\theta):=\Lb^{\circ k}\big(\Td(n+\theta)\big),
\qquad 0\le k\le n,
\]
are defined and satisfy $W_k^{(n)}(\theta)>0$ for $0\le k\le n-1$.
\end{itemize}

\begin{remark}\label{rem:hyps}
(R$'$) is precisely the quantitative reading of (R2) in Hypothesis~2.2 of
\cite{J}: positivity of $W_0,\dots,W_{n-1}$ is what it means for the $n$
real logarithms defining $D_n$ to exist; the final value $W_n=D_n$ may
well be negative (and is, e.g., for $(b,d)=(e,2)$, where
$D(0)=T_e(\Phi_{e,2}(0))<0$). See also Remarks~\ref{rem:regular} and
\ref{rem:local}.
\end{remark}

\begin{lemma}[Positivity propagation]\label{lem:positivity}
Let $T\in C^{1}$ satisfy the tetration equation $T(x+1)=a^{T(x)}$
\textup{(}$a>1$\textup{)} on its domain as in \textup{(B1)} or
\textup{(B2)}. If $T'>0$ on the single period $[0,1]$, then $T'>0$ on the
whole domain \textup{(}all of $(-2,\infty)$ in case \textup{(B1)}, all of
$[0,\infty)$ in case \textup{(B2)}\textup{)}.
\end{lemma}

\begin{proof}
Differentiating the functional equation,
\begin{equation}\label{eq:Tprimerec}
T'(x+1)=(\ln a)\,T(x+1)\,T'(x).
\end{equation}
Forward: for $x\ge0$ one has $T(x+1)>0$, so $T'>0$ on $[k,k+1]$ implies
$T'>0$ on $[k+1,k+2]$; induction covers $[0,\infty)$. Backward (case
(B1)): for $x\in[-1,0)$ we have $x+1\in[0,1)$, hence $T(x+1)>0$ and
$T'(x+1)>0$, so \eqref{eq:Tprimerec} gives
$T'(x)=T'(x+1)/\big((\ln a)T(x+1)\big)>0$; the same step with
$x\in(-2,-1)$ (where $x+1\in(-1,0)$, $T(x+1)>0$) finishes the proof.
\end{proof}

Thus (S) only needs to be checked on one fundamental interval.

\begin{remark}\label{rem:regular}
For the regular Kneser/Paulsen--Cowgill branch with real bases
$b,d>e^{1/e}$ \cite{Kneser,PC,Paulsen}, $T_a$ is real-analytic and
strictly increasing on $(-2,\infty)$ with $T_a((-2,\infty))=\R$, so
(B1)--(B2) hold. Concerning (S): we are not aware of an explicit
statement $T_a'>0$ in the literature; it is numerically evident, it is
presupposed by the numerical scheme of \cite[App.~A]{J} on its inversion
window, and by Lemma~\ref{lem:positivity} it only needs to be checked on
a single period. (By \eqref{eq:Tprimerec}, the zero set of $T_a'$ is
moreover invariant under $x\mapsto x+1$ wherever $T_a>0$, so a single
critical point would force one in every period.) Under (S), $\Td$ is
increasing, so $s_j=\Td(j)$ and
$\sup_{0\le\theta\le1}\Td(j+\theta)=\Td(j+1)$. Thus for the intended
branch the only genuinely \emph{analytic} inputs remain
Hypotheses~2.2(R2) and 2.4 of \cite{J} themselves---the same conditional
status as the phase law of \cite{J}; no new type of estimate is
introduced by Theorem~\ref{thm:D}.
\end{remark}

\begin{remark}[Local sufficiency]\label{rem:local}
Inspection of the proofs shows that of (S) only the following is actually
used: $\Td\in C^{1}$ with $\Td'>0$ on $[0,\infty)$, and $\Ab$ of class
$C^{1}$ with $\Ab'>0$ on a compact neighborhood of $D([0,1])$, together
with the Abel identity \eqref{eq:abel} near $D([0,1])$ (for
Remark~\ref{rem:ext} also near $b^{\,D([0,1])}$). The global frame (B1)
is assumed for expository convenience and holds for the intended branch
anyway.
\end{remark}

\subsection{Automatic geometric growth of the tower minima}

The derivative estimates of Section~\ref{sec:C1} pivot on the summability
$\sum_j 1/s_j<\infty$. Remarkably, this is not an additional hypothesis:
the tower functional equation supplies it.

\begin{lemma}[Automatic geometric growth]\label{lem:growth}
Assume \textup{(R1)} of Hypothesis~2.2 of \cite{J} and $s_j\to\infty$
from Hypothesis~2.4 of \cite{J}. Then
\[
s_{j+1}=d^{\,s_j}\qquad\text{for all }j\ge0,
\]
and consequently there is $J_1$ with $s_{j+1}\ge 2s_j$ for all $j\ge J_1$;
in particular
\[
\sum_{j\ge1}\frac{1}{s_j}<\infty .
\]
\end{lemma}

\begin{proof}
Fix $j$. Taking the infimum over $\theta\in[0,1]$ in
$\Td(j+1+\theta)=d^{\,\Td(j+\theta)}$ and using that $t\mapsto d^{t}$ is
increasing and continuous gives $s_{j+1}=d^{\,s_j}$; both infima run over
the same compact phase interval, and $s_j>0$ by (R1). Since $d>1$,
$d^{t}/t\to\infty$, so there is $t_0>0$ with $d^{t}\ge 2t$ for $t\ge t_0$;
choose $J_1$ with $s_j\ge t_0$ for $j\ge J_1$ (possible as
$s_j\to\infty$). Then $s_{J_1+k}\ge 2^{k}s_{J_1}$, and $\sum_j 1/s_j$
converges by comparison with a geometric series.
\end{proof}

\begin{remark}\label{rem:growthrole}
Lemma~\ref{lem:growth} shows that the qualitative divergence
$s_j\to\infty$ self-improves to geometric (indeed tower) growth. This is
special to the tetration recursion: in an abstract iterated-logarithm
scheme lacking the exact functional equation, divergence of $\sum 1/s_j$
with $\log_b\alpha\neq0$ would make the derivative products of
Section~\ref{sec:C1} degenerate to $0$ or $\infty$, so summability is
exactly the right condition---and here it is free.
\end{remark}

\subsection{Endpoint identities}
The following sharpening of Lemma~5.8 of \cite{J} will be used at both
the $C^{0}$ and the $C^{1}$ level. Only (H1) and the branch frame are
needed.

\begin{lemma}[Endpoint identities]\label{lem:endpoint}
Assume \textup{(H1)}. Then
\begin{equation}\label{eq:D1}
D(1)=b^{\,D(0)},
\qquad\text{and consequently}\qquad
\Phi_{b,d}(1)=\Phi_{b,d}(0).
\end{equation}
\end{lemma}

\begin{proof}
Since $\Lb^{\circ(n+1)}=\Lb\circ\Lb^{\circ n}$,
\[
D_{n+1}(0)=\Lb^{\circ(n+1)}\big(\Td(n+1)\big)
=\Lb\big(\Lb^{\circ n}(\Td(n+1))\big)=\Lb\big(D_n(1)\big),
\]
i.e.\ $D_n(1)=b^{\,D_{n+1}(0)}$ for every $n$ for which the quantities are
defined. Letting $n\to\infty$ and using $D_n\to D$ uniformly together with
continuity of $t\mapsto b^{t}$ gives $D(1)=b^{\,D(0)}$. Then, by the global
Abel equation \eqref{eq:abel},
\[
\Phi_{b,d}(1)=\Ab(D(1))-1=\Ab\big(b^{\,D(0)}\big)-1=\Ab(D(0))=\Phi_{b,d}(0).
\qedhere
\]
\end{proof}

In particular the $1$-periodic extension of $\Phi_{b,d}$ is continuous on
$\R$, and $x\mapsto\Phi_{b,d}(\frc{x})$ is continuous on all of $\R$,
\emph{including at the integers}. This last point is what allows
fractional parts to be composed safely below.

\section{The topological theorem}\label{sec:topo}

Throughout this section all phases are extended $1$-periodically to
$\R$; thus for a pair $(p,q)$ with package (P) we write
$\Phi_{p,q}:\R\to\R$ (continuous, $1$-periodic) and
$\Ht_{p,q}=\mathrm{id}+\Phi_{p,q}$, so that
$\Ht_{p,q}(x+1)=\Ht_{p,q}(x)+1$ and $\Phi_{p,q}(x)=\Phi_{p,q}(\frc x)$.

\begin{lemma}[Exact composition of lifted phase maps]\label{lem:comp}
Let $(p,q)$, $(q,r)$, $(p,r)$ be ordered pairs for which the package
\textup{(P)} holds. Then
\[
\Ht_{p,r}=\Ht_{p,q}\circ\Ht_{q,r}
\qquad\text{on }\R .
\]
\end{lemma}

\begin{proof}
Fix $\theta\in[0,1)$ and let $x=n+\theta$, $n\to\infty$. By
Proposition~6.1 of \cite{J}, $F_{p,r}=F_{p,q}\circ F_{q,r}$ exactly. By
(P) for $(q,r)$,
\[
y_n:=F_{q,r}(n+\theta)=n+\theta+\Phi_{q,r}(\theta)+\varepsilon_n,
\qquad \varepsilon_n\to0 .
\]
By (P) for $(p,q)$, $F_{p,q}(y)=y+\Phi_{p,q}(\frc y)+o(1)$ as
$y\to+\infty$; applying this along $y_n\to+\infty$ and using that
$u\mapsto\Phi_{p,q}(\frc u)=\Phi_{p,q}(u)$ is continuous on $\R$ (the
periodic extension; Lemma~\ref{lem:endpoint} for $(p,q)$), we get
\[
F_{p,r}(n+\theta)-(n+\theta)
=\Phi_{q,r}(\theta)+\varepsilon_n
+\Phi_{p,q}\big(\theta+\Phi_{q,r}(\theta)+\varepsilon_n\big)+o(1)
\;\longrightarrow\;
\Phi_{q,r}(\theta)+\Phi_{p,q}\big(\theta+\Phi_{q,r}(\theta)\big).
\]
The left side converges to $\Phi_{p,r}(\theta)$ by (P) for $(p,r)$. Hence
\begin{equation}\label{eq:cocycle}
\Phi_{p,r}(\theta)=\Phi_{q,r}(\theta)
+\Phi_{p,q}\big(\theta+\Phi_{q,r}(\theta)\big),
\qquad\theta\in[0,1),
\end{equation}
and by periodicity for all $\theta\in\R$. (This is Proposition~6.4 of
\cite{J}; the continuity of the periodic extension across integer values
of $\theta+\Phi_{q,r}(\theta)$, provided by Lemma~\ref{lem:endpoint}, is
what makes the limit of the composed fractional parts legitimate.)
Adding $\theta$ to \eqref{eq:cocycle},
\[
\Ht_{p,r}(\theta)
=\big(\theta+\Phi_{q,r}(\theta)\big)
+\Phi_{p,q}\big(\theta+\Phi_{q,r}(\theta)\big)
=\Ht_{p,q}\big(\Ht_{q,r}(\theta)\big).
\qedhere
\]
\end{proof}

\begin{proof}[Proof of Theorem \textup{\ref{thm:A}}]
Apply Lemma~\ref{lem:comp} to the triples $(d,b,d)$ and $(b,d,b)$; since
$\Phi_{d,d}\equiv0\equiv\Phi_{b,b}$ (diagonal packages are exact), this
gives
\[
\Ht_{d,b}\circ\Ht_{b,d}=\Ht_{d,d}=\mathrm{id},
\qquad
\Ht_{b,d}\circ\Ht_{d,b}=\Ht_{b,b}=\mathrm{id}.
\]
Thus $\Ht:=\Ht_{b,d}$ is a continuous bijection of $\R$ with continuous
inverse $\Ht_{d,b}$. A continuous injection of $\R$ is strictly
monotone; since $\Ht(x+1)=\Ht(x)+1>\Ht(x)$, it is strictly
\emph{increasing}. Because $\Ht$ commutes with $x\mapsto x+1$, it descends
to a homeomorphism $H_{b,d}$ of $\Sone$, orientation-preserving since the
lift is increasing of degree one, and $H_{b,d}^{-1}=H_{d,b}$. The groupoid
statement is Lemma~\ref{lem:comp}.
\end{proof}

\begin{remark}
Theorem~\ref{thm:A} upgrades the ``nonlinear phase cocycle'' of
\cite[Prop.~6.4]{J} to literal composition in
$\mathrm{Homeo}_{+}(\Sone)$: the family
$\{H_{p,q}\}$ is a groupoid of circle homeomorphisms, and
\eqref{eq:cocycle} is nothing but $H_{p,r}=H_{p,q}\circ H_{q,r}$ read in
the lift.
\end{remark}

\section{Antisymmetry of the mean shifts without smoothness}\label{sec:anti}

\begin{proof}[Proof of Theorem \textup{\ref{thm:B}}]
Write $\Phi=\Phi_{b,d}$, $\Psi=\Phi_{d,b}$ (both $1$-periodic and
continuous on $\R$), $\Ht=\Ht_{b,d}$, $\Ht_{d,b}=\Ht^{-1}$
(Theorem~\ref{thm:A}). The inverse case of \eqref{eq:cocycle} (triple
$(d,b,d)$) reads
\begin{equation}\label{eq:inverse}
\Psi\big(\Ht(\theta)\big)=-\Phi(\theta),\qquad\theta\in\R .
\end{equation}

Since $\Ht$ is increasing with $\Ht(\theta+1)=\Ht(\theta)+1$, the function
$\Phi=\Ht-\mathrm{id}$ is continuous and of bounded variation on every
interval of length $1$, with
$\Var_{[t,t+1]}\Phi\le\Var_{[t,t+1]}\Ht+1=\big(\Ht(t+1)-\Ht(t)\big)+1=2$.

Let $t_0:=\Ht_{d,b}(0)$, so $\Ht(t_0)=0$ and $\Ht(t_0+1)=1$. Since $\Ht$
is a continuous increasing bijection of $[t_0,t_0+1]$ onto $[0,1]$ and
$\Psi$ is continuous, the Riemann--Stieltjes change of variables
$u=\Ht(\theta)$ \cite[Thm.~7.7]{Apostol} gives, using $1$-periodicity of
$\Psi$ in the first step and \eqref{eq:inverse} in the third,
\[
\mu_{d,b}=\int_0^1\Psi(u)\,du
=\int_{t_0}^{t_0+1}\Psi\big(\Ht(\theta)\big)\,d\Ht(\theta)
=-\int_{t_0}^{t_0+1}\Phi(\theta)\,d\big(\theta+\Phi(\theta)\big)
=-\int_{t_0}^{t_0+1}\Phi\,d\theta-\int_{t_0}^{t_0+1}\Phi\,d\Phi .
\]
All Stieltjes integrals exist because the integrands are continuous and
the integrators are of bounded variation, and the splitting uses linearity
of $g\mapsto\int f\,dg$. By periodicity of $\Phi$,
$\int_{t_0}^{t_0+1}\Phi\,d\theta=\mu_{b,d}$. For the last integral,
integration by parts for Riemann--Stieltjes integrals
\cite[Thm.~7.6]{Apostol} with $f=g=\Phi$ (continuous, BV) yields
\[
2\int_{t_0}^{t_0+1}\Phi\,d\Phi
=\Phi(t_0+1)^2-\Phi(t_0)^2=0
\]
by periodicity. Hence $\mu_{d,b}=-\mu_{b,d}$.
\end{proof}

\begin{remark}\label{rem:deconditional}
The proof of Proposition~6.6 in \cite{J} is exactly this computation with
$d\Ht(\theta)=H'(\theta)\,d\theta$; the Stieltjes formulation shows that
only monotonicity---automatic by Theorem~\ref{thm:A}---is needed, not
differentiability. Together with Theorem~\ref{thm:A}, this
\emph{de-conditionalizes} Proposition~6.6 of \cite{J} completely.
\end{remark}

\begin{remark}[The smooth substitution, with the interval shift]\label{rem:smoothsub}
If in addition $\Phi,\Psi\in C^{1}$ (as under the hypotheses of
Theorem~\ref{thm:D} for both pairs, or of Theorem~\ref{thm:C}), the
substitution reads classically:
\[
\mu_{d,b}=\int_0^1\Psi(u)\,du
=\int_{\Ht(0)}^{\Ht(0)+1}\Psi(u)\,du
=\int_0^1\Psi\big(\Ht(\theta)\big)\Ht'(\theta)\,d\theta
=-\int_0^1\Phi\,(1+\Phi')\,d\theta
=-\mu_{b,d}-\tfrac12\big[\Phi^2\big]_0^1=-\mu_{b,d},
\]
where the first equality uses $1$-periodicity of $\Psi$: the map $\Ht$
carries $[0,1]$ onto $[\Ht(0),\Ht(0)+1]$ rather than onto $[0,1]$, and
periodicity supplies the missing shift. This is the computation of
\cite[Prop.~6.6]{J} with that point made explicit.
\end{remark}

\section{A minimal-regularity principle}\label{sec:minreg}

Theorem~\ref{thm:C} is a purely formal statement about pairs of
$1$-periodic $C^{1}$ functions coupled by the inverse phase equations;
it consumes no channel estimates and no derivative formulas.

\begin{proof}[Proof of Theorem \textup{\ref{thm:C}}]
The equations \eqref{eq:invpair} give, for every $x\in\R$,
\[
\Kt\big(\Ht(x)\big)=\Ht(x)+\Psi\big(\Ht(x)\big)
=x+\Phi(x)-\Phi(x)=x ,
\]
and symmetrically $\Ht(\Kt(u))=u$; thus $\Kt\circ\Ht=\Ht\circ\Kt=
\mathrm{id}$ with $\Ht,\Kt\in C^{1}(\R)$. Differentiating the first
identity by the chain rule,
\[
\Kt'\big(\Ht(x)\big)\,\Ht'(x)=1\qquad\text{for all }x\in\R ,
\]
so $\Ht'$ has no zeros. Being continuous on the connected set $\R$,
$\Ht'$ has a constant sign, and
\[
\int_0^1\Ht'(x)\,dx=\Ht(1)-\Ht(0)=1>0
\]
($1$-periodicity of $\Phi$) forces the sign to be positive: $\Ht'>0$
everywhere, and $\Kt'=1/(\Ht'\circ\Kt)>0$. A $C^{1}$ map of $\R$ with
everywhere positive derivative and $\Ht(x+1)=\Ht(x)+1$ is a strictly
increasing bijection of $\R$; its inverse $\Kt$ is $C^{1}$, both commute
with the unit translation, and so they descend to mutually inverse
orientation-preserving $C^{1}$ diffeomorphisms of $\Sone$. Finally, under
(H1)+(H2) the equations \eqref{eq:invpair} for $\Phi=\Phi_{b,d}$,
$\Psi=\Phi_{d,b}$ are the cases $(d,b,d)$ and $(b,d,b)$ of
Lemma~\ref{lem:comp}.
\end{proof}

\begin{example}[One-sided $C^{1}$ regularity does not suffice]\label{ex:counter}
Let
\[
\Phi_0(x):=-\frac{\sin(2\pi x)}{2\pi},
\qquad
\Ht_0:=\mathrm{id}+\Phi_0 .
\]
Then $\Phi_0$ is $1$-periodic and real-analytic,
$\Ht_0'(x)=1-\cos(2\pi x)\ge0$ vanishes exactly at $x\in\Z$, and $\Ht_0$
is strictly increasing: for $y>x$,
$\Ht_0(y)-\Ht_0(x)=\int_x^y(1-\cos2\pi t)\,dt>0$ since the integrand
vanishes only at isolated points. Hence $\Ht_0$ is the lift of a
degree-one $C^{1}$ circle \emph{homeomorphism} which is \emph{not} a
$C^{1}$ diffeomorphism: near $x=0$,
\[
\Ht_0(x)=\frac{2\pi^{2}}{3}\,x^{3}+O(x^{5}),
\]
so the inverse behaves like a cube root at $0$ and is not differentiable
there; the inverse phase $\Psi_0=\Ht_0^{-1}-\mathrm{id}$ is continuous
but not $C^{1}$. Consequently, $C^{1}$ regularity plus topological
invertibility of the \emph{forward} lift alone do not imply the
diffeomorphism property. Theorem~\ref{thm:C} genuinely consumes $C^{1}$
regularity of both phases, while Theorem~\ref{thm:D} supplies the
positivity $\Ht'>0$ for the forward pair outright.
\end{example}

\begin{corollary}\label{cor:redundant}
Assume \textup{(H1)} and \textup{(H2)}. If $\Phi_{b,d}$ and $\Phi_{d,b}$
are of class $C^{1}$ on the circle---however this is established---then
$H_{b,d}$ is an orientation-preserving $C^{1}$ circle diffeomorphism with
$C^{1}$ inverse $H_{d,b}$. In particular the separate diffeomorphism
assumption in \cite[Prop.~6.6]{J} is redundant given two-sided $C^{1}$
phase limits.
\end{corollary}

\begin{remark}[Logical role]\label{rem:role}
Theorem~\ref{thm:C} decouples the diffeomorphism property from the
specific analytic mechanism: any route to two-sided $C^{1}$ phases
inherits the full conclusion. For instance, a proof of strip analyticity
(\cite[Conj.~8.1]{J}; Remark~\ref{rem:Cr}) for both ordered pairs would
immediately produce mutually inverse real-analytic circle
diffeomorphisms. Within the present note it is not strictly needed,
since Theorem~\ref{thm:D}(i) proves $\Ht'>0$ directly.
\end{remark}

\section{The analytic $C^{1}$ theorem}\label{sec:C1}

Standing assumptions for this section: (H1), (S), (R$'$), together
with the branch frame (B1)--(B2). We abbreviate
$W_k=W_k^{(n)}(\theta)=\Lb^{\circ k}(\Td(n+\theta))$, so
$W_0=\Td(n+\theta)$, $W_{k+1}=\Lb(W_k)$, $W_n=D_n(\theta)$, and, by
(R$'$), $W_k>0$ for $0\le k\le n-1$ and $n\ge N_0$. Since
$W_{k+1}=\Lb(W_k)$ with $W_k>0$ is equivalent to $W_k=b^{\,W_{k+1}}$, we
record the inversion
\begin{equation}\label{eq:inversion}
W_{n-j}^{(n)}(\theta)=\expb^{\circ j}\big(D_n(\theta)\big),
\qquad 1\le j\le n .
\end{equation}

\subsection{The exact derivative product identity}

\begin{lemma}[Exact product identity]\label{lem:product}
Let $n\ge\max(N_0,2)$ and $\theta\in[0,1]$. Then $D_n\in C^{1}([0,1])$,
and with
\begin{equation}\label{eq:sigmadef}
1+\sigma_j^{(n)}(\theta)
:=\frac{W_{n-j}^{(n)}(\theta)}{\alpha\,\Td(j+\theta)}
=\frac{\expb^{\circ j}\big(D_n(\theta)\big)}{\alpha\,\Td(j+\theta)}\;>\;0,
\qquad 1\le j\le n-1,
\end{equation}
one has the \emph{exact} identity
\begin{equation}\label{eq:exact}
D_n'(\theta)
=\frac{\alpha\,\Td'(\theta)}{\displaystyle\prod_{j=1}^{n-1}\big(1+\sigma_j^{(n)}(\theta)\big)} ,
\qquad\text{and}\qquad
\sigma_{n-1}^{(n)}\equiv 0 .
\end{equation}
\end{lemma}

\begin{proof}
By (S) and (R$'$), each map in the chain
$\theta\mapsto\Td(n+\theta)=W_0\mapsto W_1\mapsto\cdots\mapsto W_n=D_n(\theta)$
is $C^{1}$ along the orbit (the logarithms are applied to positive
quantities, which range over a compact subset of $(0,\infty)$ as
$\theta\in[0,1]$), so $D_n\in C^{1}([0,1])$ and the chain rule with
$\Lb'(y)=1/(y\ln b)$ gives
\begin{equation}\label{eq:chain}
D_n'(\theta)=\frac{\Td'(n+\theta)}{(\ln b)^{n}\,\prod_{k=0}^{n-1}W_k}.
\end{equation}
Differentiating $\Td(x+1)=d^{\Td(x)}$ yields
$\Td'(x+1)=(\ln d)\Td(x+1)\Td'(x)$, and iterating from $x=\theta$,
\begin{equation}\label{eq:Tdprime}
\Td'(n+\theta)=(\ln d)^{n}\,\Td'(\theta)\prod_{i=1}^{n}\Td(i+\theta).
\end{equation}
Next, the first peel is exact:
\[
W_1=\Lb\big(\Td(n+\theta)\big)=\Lb\big(d^{\,\Td(n-1+\theta)}\big)
=\alpha\,\Td(n-1+\theta),
\]
which is \eqref{eq:sigmadef} with $j=n-1$ and $\sigma_{n-1}^{(n)}=0$.
Using \eqref{eq:sigmadef} for $1\le k\le n-1$ (i.e.\ $j=n-k$),
\[
\prod_{k=0}^{n-1}W_k
=\Td(n+\theta)\cdot\alpha^{\,n-1}\prod_{i=1}^{n-1}\Td(i+\theta)
\cdot\prod_{j=1}^{n-1}\big(1+\sigma_j^{(n)}(\theta)\big).
\]
Substituting this and \eqref{eq:Tdprime} into \eqref{eq:chain}, the factor
$(\ln d)^{n}/(\ln b)^{n}=\alpha^{n}$ appears, the products
$\prod_{i=1}^{n}\Td(i+\theta)$ and
$\Td(n+\theta)\prod_{i=1}^{n-1}\Td(i+\theta)$ cancel exactly, and
$\alpha^{n}/\alpha^{n-1}=\alpha$ survives:
\[
D_n'(\theta)
=\frac{\alpha^{n}\,\Td'(\theta)\prod_{i=1}^{n}\Td(i+\theta)}
{\alpha^{\,n-1}\,\Td(n+\theta)\prod_{i=1}^{n-1}\Td(i+\theta)\,
\prod_{j=1}^{n-1}(1+\sigma_j^{(n)})}
=\frac{\alpha\,\Td'(\theta)}{\prod_{j=1}^{n-1}(1+\sigma_j^{(n)})}.
\]
The identification \eqref{eq:sigmadef} of $W_{n-j}$ with
$\expb^{\circ j}(D_n)$ is \eqref{eq:inversion}.
\end{proof}

Note that \eqref{eq:exact} is an algebraic identity: it holds for
\emph{any} $C^{1}$ germ propagated by the tower recursion, independently
of the branch choice; see Remark~\ref{rem:numerics} for a high-precision
numerical confirmation.

\begin{remark}[Ratio identity and an alternative route]\label{rem:ratio}
Dividing the chain-rule products for consecutive tails yields a companion
identity. By \cite[eq.~(22)]{J}, $D_{n+1}(\theta)=\Lb^{\circ n}\big(\alpha
\Td(n+\theta)\big)$; writing $\widetilde V_k:=\Lb^{\circ k}(\alpha
\Td(n+\theta))$ and $\delta_k:=\widetilde V_k-W_k$---the damped
differences of \cite[Lemma~5.3]{J}, with $\delta_1=\Lb(\alpha)$---one
obtains, since $\widetilde V_0=\alpha W_0$ and all quantities are positive
by \textup{(R$'$)} at level $n+1$,
\[
\frac{D_{n+1}'(\theta)}{D_n'(\theta)}
=\prod_{k=1}^{n-1}\frac{W_k}{W_k+\delta_k}.
\]
The huge factor $\Td'(n+\theta)$ cancels identically. With the estimates
of \cite[Lemma~5.3]{J} one gets
$\|\log D_{n+1}'-\log D_n'\|_\infty=O(1/s_{n-1})$, summable in $n$ by
Lemma~\ref{lem:growth}; hence $(\log D_n')$ is uniformly Cauchy, giving
an alternative proof of Proposition~\ref{prop:C1conv} that inherits the
hyper-exponential speed of \cite[Lemma~5.4]{J}. We follow the closed-form
route below because it produces the limit formula \eqref{eq:Dprime} and
the positivity bound in the same stroke.
\end{remark}

\subsection{Tail estimates}

\begin{lemma}[Recursion and tail bound]\label{lem:tail}
For $n\ge\max(N_0,2)$, $\theta\in[0,1]$, and $1\le j\le n-2$,
\begin{equation}\label{eq:rec}
\sigma_j^{(n)}(\theta)
=\frac{\Lb(\alpha)+\Lb\big(1+\sigma_{j+1}^{(n)}(\theta)\big)}
{\alpha\,\Td(j+\theta)} ,
\qquad \sigma_{n-1}^{(n)}=0 .
\end{equation}
Let
\[
C_0:=\frac{|\Lb(\alpha)|+\dfrac{\ln 2}{\ln b}}{\alpha},
\qquad
J_0:=\min\Big\{J:\ \frac{C_0}{s_j}\le\tfrac12\ \text{for all }j\ge J\Big\}
\]
\textup{(}finite since $s_j\to\infty$\textup{)}. Then for all
$n\ge\max(N_0,2)$ and $\theta\in[0,1]$,
\begin{equation}\label{eq:tailbound}
\big|\sigma_j^{(n)}(\theta)\big|\le\frac{C_0}{s_j}\le\frac12
\qquad\text{for } J_0\le j\le n-1 .
\end{equation}
\end{lemma}

\begin{proof}
For $1\le j\le n-2$, using \eqref{eq:sigmadef} at level $j+1$, the
identity $\Lb(\Td(j+1+\theta))=\Lb(d^{\Td(j+\theta)})=\alpha\Td(j+\theta)$,
and additivity of $\Lb$ on positive factors,
\[
W_{n-j}=\Lb(W_{n-j-1})
=\Lb\Big(\alpha\,\Td(j+1+\theta)\big(1+\sigma_{j+1}^{(n)}\big)\Big)
=\Lb(\alpha)+\alpha\,\Td(j+\theta)+\Lb\big(1+\sigma_{j+1}^{(n)}\big);
\]
dividing by $\alpha\Td(j+\theta)$ gives \eqref{eq:rec}. For
\eqref{eq:tailbound} we induct downward on $j$ from $j=n-1$ to $j=J_0$.
The base case is $\sigma_{n-1}^{(n)}=0$. If
$|\sigma_{j+1}^{(n)}|\le\frac12$ (with $J_0\le j\le n-2$), then
$|\ln(1+\sigma_{j+1}^{(n)})|\le\ln2$, so \eqref{eq:rec} gives
\[
\big|\sigma_j^{(n)}\big|
\le\frac{|\Lb(\alpha)|+\frac{\ln2}{\ln b}}{\alpha\,\Td(j+\theta)}
\le\frac{C_0}{s_j}\le\frac12 .
\qedhere
\]
\end{proof}

\begin{lemma}[Head control and limit factors]\label{lem:head}
There is $N_1\ge\max(N_0,2)$ such that for every fixed $j\ge1$:
\begin{enumerate}
\item[\textup{(a)}] $1+\sigma_j^{(n)}\to 1+\sigma_j$ uniformly on
$[0,1]$ as $n\to\infty$, where
\[
1+\sigma_j(\theta):=\frac{E_j(\theta)}{\alpha\,\Td(j+\theta)},
\qquad
E_j(\theta):=\expb^{\circ j}\big(D(\theta)\big)>0 ;
\]
\item[\textup{(b)}] there are constants $0<m_j\le M_j<\infty$ with
$m_j\le 1+\sigma_j^{(n)}(\theta)\le M_j$ for all $n\ge N_1$,
$\theta\in[0,1]$, and $m_j\le 1+\sigma_j(\theta)\le M_j$;
\item[\textup{(c)}] $|\sigma_j(\theta)|\le C_0/s_j$ for $j\ge J_0$.
\end{enumerate}
\end{lemma}

\begin{proof}
By (H1), $D_n\to D$ uniformly on $[0,1]$; choose $N_1\ge\max(N_0,2)$ with
$\|D_n-D\|_\infty\le1$ for $n\ge N_1$ and set
$K:=[\min_{[0,1]}D-1,\ \max_{[0,1]}D+1]$, a compact interval containing
all values $D_n(\theta)$ ($n\ge N_1$) and $D(\theta)$. Each
$\expb^{\circ j}$ is $C^{1}$, hence Lipschitz on $K$, and
$\expb^{\circ j}(K)=:[c_j,C_j]\subset(0,\infty)$ for $j\ge1$. By
\eqref{eq:inversion}, $W_{n-j}^{(n)}=\expb^{\circ j}(D_n)\to
\expb^{\circ j}(D)=E_j$ uniformly, which gives (a); moreover
$E_j\ge c_j>0$. By (B2) and continuity, $\Td(j+\cdot)$ has a finite
supremum $S_j$ and infimum $s_j>0$ on $[0,1]$, so (b) holds with
$m_j=c_j/(\alpha S_j)$, $M_j=C_j/(\alpha s_j)$. Passing to the limit
$n\to\infty$ in \eqref{eq:tailbound} gives (c).
\end{proof}

\subsection{Uniform convergence of the product and $C^{1}$ convergence}

\begin{lemma}[Uniform product convergence]\label{lem:prodconv}
Set
\[
\Pi_n(\theta):=\prod_{j=1}^{n-1}\big(1+\sigma_j^{(n)}(\theta)\big),
\qquad
\Pi(\theta):=\prod_{j=1}^{\infty}\big(1+\sigma_j(\theta)\big).
\]
Then $\Pi$ converges absolutely and uniformly on $[0,1]$, is continuous,
satisfies $0<c_\Pi\le\Pi\le C_\Pi<\infty$, and $\Pi_n\to\Pi$ uniformly on
$[0,1]$.
\end{lemma}

\begin{proof}
All factors are positive (Lemma~\ref{lem:product},
Lemma~\ref{lem:head}(b)), so we may work with logarithms. For
$|s|\le\frac12$ one has $|\ln(1+s)|\le 2|s|$; hence by
\eqref{eq:tailbound} and Lemma~\ref{lem:head}(c),
\begin{equation}\label{eq:logtail}
\big|\ln\big(1+\sigma_j^{(n)}\big)\big|\le\frac{2C_0}{s_j},
\quad
\big|\ln\big(1+\sigma_j\big)\big|\le\frac{2C_0}{s_j}
\qquad (j\ge J_0),
\end{equation}
uniformly in $n$ and $\theta$. By Lemma~\ref{lem:growth},
$\sum_j 2C_0/s_j<\infty$, which already gives absolute and uniform
convergence of $\ln\Pi=\sum_{j\ge1}\ln(1+\sigma_j)$ (the finitely many
terms $j<J_0$ are bounded by Lemma~\ref{lem:head}(b)), hence continuity
of $\Pi$ and the two-sided bound with
$\ln C_\Pi=-\ln c_\Pi:=\sum_{j<J_0}\max(|\ln m_j|,|\ln M_j|)
+\sum_{j\ge J_0}2C_0/s_j$.

For $\Pi_n\to\Pi$: given $\varepsilon>0$, choose $J\ge J_0$ with
$\sum_{j\ge J}2C_0/s_j<\varepsilon/3$ (Lemma~\ref{lem:growth}). For
$n-1\ge J$,
\[
\big|\ln\Pi_n-\ln\Pi\big|
\le\sum_{j=1}^{J-1}\Big|\ln\big(1+\sigma_j^{(n)}\big)-\ln\big(1+\sigma_j\big)\Big|
+\sum_{j=J}^{n-1}\big|\ln\big(1+\sigma_j^{(n)}\big)\big|
+\sum_{j=J}^{\infty}\big|\ln\big(1+\sigma_j\big)\big| .
\]
The two tails are each $<\varepsilon/3$ by \eqref{eq:logtail}. In the
head, each of the finitely many terms tends to $0$ uniformly by
Lemma~\ref{lem:head}(a) and the Lipschitz continuity of $\ln$ on
$[m_j\wedge\inf(1+\sigma_j),\infty)\subset[m_j,\infty)$. Hence
$\limsup_n\|\ln\Pi_n-\ln\Pi\|_\infty\le\tfrac{2\varepsilon}{3}<\varepsilon$,
so $\ln\Pi_n\to\ln\Pi$ uniformly, and applying $\exp$ (uniformly
continuous on the compact range) gives $\Pi_n\to\Pi$ uniformly.
\end{proof}

\begin{proposition}[$C^{1}$ convergence of the tails]\label{prop:C1conv}
$D_n\to D$ in $C^{1}([0,1])$; $D\in C^{1}([0,1])$; and
\[
D'(\theta)=\frac{\alpha\,\Td'(\theta)}{\Pi(\theta)}
=\alpha\,\Td'(\theta)\prod_{j=1}^{\infty}
\frac{\alpha\,\Td(j+\theta)}{E_j(\theta)}
\;\ge\;\frac{\alpha\,\min_{[0,1]}\Td'}{C_\Pi}\;>\;0 .
\]
\end{proposition}

\begin{proof}
By Lemma~\ref{lem:product} and Lemma~\ref{lem:prodconv},
$D_n'=\alpha\Td'/\Pi_n\to\alpha\Td'/\Pi=:G$ uniformly on $[0,1]$, with $G$
continuous and, by (S) and Lemma~\ref{lem:prodconv}, bounded below by the
stated positive constant. Since also $D_n\to D$ uniformly (H1), the
standard theorem on differentiation of uniform limits
\cite[Thm.~7.17]{Rudin} yields $D\in C^{1}$ and $D'=G$. The product form
of $G$ is \eqref{eq:sigmadef}--\eqref{eq:exact} in the limit, i.e.\
$(1+\sigma_j)^{-1}=\alpha\Td(j+\theta)/E_j$.
\end{proof}

\subsection{Seam matching}

\begin{lemma}[$C^{1}$ seam identities]\label{lem:seam}
Under the standing assumptions,
\[
D(1)=b^{\,D(0)},\qquad
D'(1)=(\ln b)\,D(1)\,D'(0),
\]
and consequently, with $\Phi=\Phi_{b,d}=\Ab\circ D-\mathrm{id}$ on $[0,1]$,
\[
\Phi(1)=\Phi(0),
\qquad
\Phi'(1)=\Phi'(0).
\]
\end{lemma}

\begin{proof}
The value identities are Lemma~\ref{lem:endpoint}. For the derivative
identity, evaluate the product formula of
Proposition~\ref{prop:C1conv} at $\theta=1$ and $\theta=0$. First,
$\Td'(1)=(\ln d)\Td(1)\Td'(0)$ by \eqref{eq:Tprimerec}. Second, by
$D(1)=b^{D(0)}$,
\[
E_j(1)=\expb^{\circ j}\big(D(1)\big)
=\expb^{\circ j}\big(b^{\,D(0)}\big)
=\expb^{\circ(j+1)}\big(D(0)\big)
=E_{j+1}(0),
\qquad j\ge 1 .
\]
Hence, shifting the product index ($i=j+1$),
\[
D'(1)=\alpha\,(\ln d)\Td(1)\Td'(0)
\prod_{j=1}^{\infty}\frac{\alpha\,\Td(j+1)}{E_{j+1}(0)}
=\alpha\,(\ln d)\Td(1)\Td'(0)
\prod_{i=2}^{\infty}\frac{\alpha\,\Td(i)}{E_{i}(0)} ,
\]
while $D'(0)=\alpha\,\Td'(0)\prod_{i=1}^{\infty}\alpha\Td(i)/E_i(0)$.
All products converge to nonzero limits (Lemma~\ref{lem:prodconv}), so
the termwise quotient is legitimate:
\[
\frac{D'(1)}{D'(0)}
=(\ln d)\,\Td(1)\cdot\Big(\frac{\alpha\,\Td(1)}{E_1(0)}\Big)^{-1}
=\frac{\ln d}{\alpha}\,E_1(0)
=(\ln b)\,b^{\,D(0)}
=(\ln b)\,D(1).
\]
Now differentiate the global Abel equation \eqref{eq:abel} (both sides
$C^{1}$ by (S)) at $y=D(0)$:
\[
\Ab'\big(b^{\,y}\big)\,b^{\,y}\ln b=\Ab'(y)
\quad\Longrightarrow\quad
\Ab'(D(1))\,(\ln b)\,D(1)=\Ab'(D(0)) .
\]
Therefore
\[
\Phi'(1)+1=\Ab'(D(1))\,D'(1)
=\frac{\Ab'(D(0))}{(\ln b)D(1)}\cdot(\ln b)D(1)\,D'(0)
=\Ab'(D(0))\,D'(0)=\Phi'(0)+1 .
\qedhere
\]
\end{proof}

\begin{remark}[Seam via extended phases]\label{rem:ext}
Alternatively: the exact transport
$D_n(1+\theta)=b^{\,D_{n+1}(\theta)}$ (the inversion \eqref{eq:inversion}
one level up; positivity by \textup{(R$'$)} at level $n{+}1$) carries the
uniform convergence of $D_{n+1}$ on $[0,1]$ to that of $D_n$ on $[1,2]$,
and every estimate of this section holds verbatim with phase range
$[0,2]$ (with $s_j$ replaced by $\min(s_j,s_{j+1})$ in the tail bounds).
Running Proposition~\ref{prop:C1conv} on $[0,2]$ gives
$D\in C^{1}([0,2])$ with $D(1+\theta)=b^{\,D(\theta)}$, whence
$\Phi(1+\theta)=\Phi(\theta)$ holds as an identity of $C^{1}$ functions on
$[0,1]$---including the derivative match at the seam---without touching
the product formula. This formulation extends verbatim to higher
derivatives; see Remark~\ref{rem:Cr}.
\end{remark}

\subsection{Proof of Theorem \ref{thm:D} and corollaries}

\begin{proof}[Proof of Theorem \textup{\ref{thm:D}}]
\textup{(i)} is Proposition~\ref{prop:C1conv}.

\textup{(ii)} By (S) and \eqref{eq:Abprime}, $\Ab\in C^{1}(\R)$ with
$\Ab'>0$, so $\Phi=\Ab\circ D-\mathrm{id}\in C^{1}([0,1])$ with
$\Phi'=(\Ab'\circ D)D'-1$. By Lemma~\ref{lem:seam}, $\Phi$ has matching
values and one-sided derivatives at the endpoints, so its $1$-periodic
extension is $C^{1}$ on $\R$ (the extension is differentiable at the
integers with derivative $\Phi'(0)=\Phi'(1)$, and $\Phi'$ extends
continuously and periodically). Hence $\Ht=\Ht_{b,d}=\mathrm{id}+\Phi\in
C^{1}(\R)$, $\Ht(x+1)=\Ht(x)+1$, and on $[0,1]$ (hence on $\R$ by
periodicity of $\Ht'$)
\[
\Ht'=1+\Phi'=(\Ab'\circ D)\,D'
\;\ge\;\delta_0:=\min_{[0,1]}\big[(\Ab'\circ D)\,D'\big]>0 ,
\]
using Proposition~\ref{prop:C1conv} and continuity on the compact
interval. For the closed form \eqref{eq:Hprime}: since
$\Ht(\theta)=\Ab(D(\theta))$, we have $\Tb(\Ht(\theta))=D(\theta)$, hence
by iterating $\Tb(x+1)=b^{\Tb(x)}$,
\[
\Tb\big(j+\Ht(\theta)\big)=\expb^{\circ j}\big(\Tb(\Ht(\theta))\big)
=\expb^{\circ j}\big(D(\theta)\big)=E_j(\theta),
\qquad j\ge1,
\]
and $\Ab'(D(\theta))=1/\Tb'(\Ab(D(\theta)))=1/\Tb'(\Ht(\theta))$ by
\eqref{eq:Abprime}; substituting into
$\Ht'=(\Ab'\circ D)D'$ and \eqref{eq:Dprime} gives \eqref{eq:Hprime}.

\textup{(iii)} $\Ht$ is $C^{1}$ with $\Ht'\ge\delta_0>0$ and
$\Ht(x+1)=\Ht(x)+1$; hence $\Ht$ is a strictly increasing $C^{1}$
bijection of $\R$ (surjectivity from the translation property), its
inverse is $C^{1}$ by the inverse function theorem, and both $\Ht$ and
$\Ht^{-1}$ commute with the unit translation. Therefore $\Ht$ descends to
an orientation-preserving $C^{1}$ diffeomorphism $H_{b,d}$ of $\Sone$.
This is precisely the hypothesis of Proposition~6.6 of \cite{J}. If
moreover (H2) holds, Theorem~\ref{thm:A} identifies
$\Ht^{-1}=\Ht_{d,b}=\mathrm{id}+\Phi_{d,b}$, so
$\Phi_{d,b}=\Ht^{-1}-\mathrm{id}\in C^{1}$, i.e.\ the reversed phase is
automatically $C^{1}$ on the circle. (Alternatively, once both phases
are known to be $C^{1}$, Theorem~\ref{thm:C} applies.)

\textup{(iv)} By Proposition~5.1 of \cite{J},
$R_n(\theta)=F_{b,d}(n+\theta)-(n+\theta)=\Ab(D_n(\theta))-\theta$
exactly. Hence $R_n\in C^{1}([0,1])$ and
$R_n'=(\Ab'\circ D_n)\,D_n'-1$. Since $D_n\to D$ uniformly with values in
the compact $K$ (Lemma~\ref{lem:head}), $\Ab'$ is uniformly continuous on
$K$, and $D_n'\to D'$ uniformly (Proposition~\ref{prop:C1conv}), we get
$R_n'\to(\Ab'\circ D)D'-1=\Phi'$ uniformly; together with $R_n\to\Phi$
uniformly (Theorem~5.7 of \cite{J}) this is $C^{1}$ convergence. Finally
$F_{b,d}'(n+\theta)=1+R_n'(\theta)\to 1+\Phi'(\theta)=\Ht'(\theta)$.
\end{proof}

\begin{corollary}[$C^{1}$ groupoid]\label{cor:groupoid}
If the hypotheses of Theorem~\textup{\ref{thm:D}} hold for all ordered
pairs from $\{b,c,d\}$, then all maps $H_{p,q}$ are orientation-preserving
$C^{1}$ circle diffeomorphisms and
$H_{p,r}=H_{p,q}\circ H_{q,r}$ in $\mathrm{Diff}^{1}_{+}(\Sone)$; the
identity $\Phi_{p,r}=\Phi_{q,r}+\Phi_{p,q}\circ(\mathrm{id}+\Phi_{q,r})$
of \cite[Prop.~6.4]{J} is its expression in lifts.
\end{corollary}

\begin{proof}
Combine Theorem~\ref{thm:D} with Lemma~\ref{lem:comp}.
\end{proof}

\begin{corollary}[Proposition~6.6 of \cite{J}, discharged]\label{cor:66}
Under \textup{(H1)} and \textup{(H2)}: $\mu_{b,d}=-\mu_{d,b}$
(Theorem~\ref{thm:B}; no smoothness). Under \textup{(H1)},
\textup{(S)}, \textup{(R$'$)}: $H_{b,d}$ is an orientation-preserving
$C^{1}$ circle diffeomorphism, i.e.\ the literal hypothesis of
\cite[Prop.~6.6]{J} holds; if \textup{(H2)} holds as well, so does its
inverse-side counterpart, with $H_{d,b}=H_{b,d}^{-1}$ of class $C^{1}$.
\end{corollary}

\begin{proof}
Immediate from Theorems~\ref{thm:B} and \ref{thm:D}.
\end{proof}

\section{Remarks}\label{sec:remarks}

\begin{remark}[The mismatch sits one level down]\label{rem:onelevel}
By \eqref{eq:sigmadef} in the limit and Lemma~\ref{lem:head}(c),
\[
\Tb\big(j+\Ht_{b,d}(\theta)\big)=E_j(\theta)
=\alpha\,\Td(j+\theta)\,\big(1+\sigma_j(\theta)\big),
\qquad
|\sigma_j(\theta)|\le\frac{C_0}{s_j}\quad(j\ge J_0),
\]
uniformly in $\theta$. Thus measuring the $d$-tower at height $j+\theta$
in $b$-coordinates at the \emph{phase-corrected} height
$j+\Ht_{b,d}(\theta)$ reproduces the value up to the single factor
$\alpha$ and a hyper-exponentially small relative error: the base
mismatch survives exactly one tower level down. This quantifies, at the
level of values, the damping mechanism behind Theorem~5.6 of \cite{J},
and it is the reason the derivative product \eqref{eq:Hprime} converges
so violently fast.
\end{remark}

\begin{remark}[Factorization $D=\Tb\circ\Ht$]\label{rem:factor}
Since $\Ht(\theta)=\Ab(D(\theta))$, one has $D=\Tb\circ\Ht$ on $[0,1]$:
the limiting tail is the phase map read in the $\Tb$-coordinate. With the
extension of Remark~\ref{rem:ext}, the translation identity
$D(1+\theta)=b^{\,D(\theta)}$ is precisely the tetration equation of
$\Tb$ transported by $\Ht$. Given the branch regularity (S), regularity
of $D$ and regularity of $\Ht$ are therefore equivalent.
\end{remark}

\begin{remark}[$C^{r}$ and analytic upgrades]\label{rem:Cr}
If in (S) the branches are of class $C^{r}$ ($r\ge2$), the same scheme
yields $\Phi_{b,d}\in C^{r}(\Sone)$ and a $C^{r}$ diffeomorphism:
differentiate
$\log D_n'=\log(\alpha\Td')-\sum_j\log(1+\sigma_j^{(n)})$ repeatedly, or
run the extension bookkeeping of Remark~\ref{rem:ext}; the tail control
is again Lemma~\ref{lem:tail} together with Lemma~\ref{lem:growth}. The
proof is likewise stable under complexification: if the channel
estimates and the uniform convergence $D_n\to D$ hold on a strip
$\{|\mathrm{Im}\,\theta|<\sigma_0\}$ (with $\Td$, $\Tb$ holomorphic
there, as they are for the regular branch), Weierstrass' theorem upgrades
Proposition~\ref{prop:C1conv} to locally uniform convergence of
holomorphic functions; then $D$, $\Phi_{b,d}$ and $\Ht_{b,d}$ are
real-analytic, the Fourier coefficients of $P_{b,d}$ decay exponentially,
and Conjecture~8.1 of \cite{J} follows---by Theorem~\ref{thm:C}, both
phase maps are then automatically mutually inverse \emph{analytic} circle
diffeomorphisms. The product identity \eqref{eq:exact} is already the
correct analytic object for that program; we do not pursue the complex
channel estimates here.
\end{remark}

\begin{remark}[What remains conditional]\label{rem:conditional}
Theorems~\ref{thm:A} and \ref{thm:B} consume only the conclusions of the
phase theorem of \cite{J} (for both ordered pairs); Theorem~\ref{thm:C}
is unconditional as a statement about $C^{1}$ data; and
Theorem~\ref{thm:D} consumes, beyond (H1), only (S) (true for the
regular branch, and reducible to one period by
Lemma~\ref{lem:positivity}) and (R$'$) (the quantitative reading of the
paper's (R2)); the growth input is supplied by Lemma~\ref{lem:growth}.
Thus the results are conditional in exactly the same sense as the main
theorem of \cite{J}: everything reduces to the real-channel and
peeling-growth estimates for the regular branch (Section~8.3 of
\cite{J}); no additional analytic difficulty is introduced by the
diffeomorphism statement.
\end{remark}

\begin{remark}[Numerical validation]\label{rem:numerics}
\textup{(i)} \emph{The exact identity.} Identity \eqref{eq:exact} holds
for any $C^{1}$ germ propagated by the tower recursion, so it can be
tested independently of any branch construction: take an arbitrary
$C^{1}$ function $f>0$ near $\theta_0$, define $T_0=f$,
$T_j=d^{\,T_{j-1}}$, $D_n=\Lb^{\circ n}(T_n)$, and compare
$D_n'(\theta_0)$ (central finite differences, step $10^{-30}$, $60$-digit
arithmetic) with
$\alpha f'(\theta_0)\prod_{j=1}^{n-1}\alpha T_j(\theta_0)/
\expb^{\circ j}(D_n(\theta_0))$. For $(b,d)=(e,2)$ at $n=4,5$ and
$(b,d)=(2,10)$ at $n=3,4$ (seeds $1+\theta+0.3\sin(2.7\theta+0.4)$ and
$0.2(1+\theta)+0.05\sin(2.7\theta+0.4)$, $\theta_0=0.37$), the two sides
agree to $31$--$33$ significant digits---the finite-difference accuracy
limit---and the inversion \eqref{eq:inversion} holds to working
precision ($<10^{-58}$ relative).

\textup{(ii)} \emph{End-to-end validation on explicit $C^{1}$ branches.}
For a base $a>1$ let
\[
q_a(x):=1+c_1x+c_2x^{2}\ \text{ on }[-1,0],
\qquad
c_1=\frac{2\ln a}{1+\ln a},\quad c_2=c_1-1,
\]
extended to $(-2,\infty)$ by the tetration equation. Then $T_a$ is a
genuine $C^{1}$ solution of the frame (B1)--(B2) with $T_a'>0$: the
$C^{1}$ match at the integers is exactly the constraint defining $c_1$,
one has $q_a'\ge\min(c_1,2-c_1)>0$ on $[-1,0]$, and
Lemma~\ref{lem:positivity} propagates; for $a=e$ this is the classical
construction with linear critical piece, $T_e(\theta)=e^{\theta}$ on
$[0,1]$. On these branches all hypotheses of Theorem~\ref{thm:D} are
directly verifiable, and the recursion \eqref{eq:rec} doubles as the
numerically stable evaluation scheme---the interleaved peeling of
\cite[App.~A]{J} in $\sigma$-coordinates: truncating $\sigma_{j^*}=0$ at
the first level with $\Td(j^*+\theta)$ beyond working precision
reproduces the finite-$n$ objects exactly and the limits to full
precision. At $60$-digit precision, for the pairs $(e,2)$ and $(2,e)$:
the $\sigma$-scheme agrees with direct tower-plus-peeling evaluation to
below $10^{-60}$; the increments $\|\Phi_{n+1}-\Phi_n\|$ at
$\theta=0.37$ fall as $2.2\cdot10^{-3}$ $(n{=}4{\to}5)$ and
$4.6\cdot10^{-19}$ $(n{=}5{\to}6)$ for $(e,2)$, and as $8.5\cdot10^{-4}$
$(n{=}3{\to}4)$ and $2.4\cdot10^{-34}$ $(n{=}4{\to}5)$ for $(2,e)$,
exhibiting the predicted hyper-exponential collapse; the closed formula
\eqref{eq:Hprime} matches central finite differences at the
finite-difference accuracy limit ($\approx10^{-36}$); the seam residuals
$\Phi(1)-\Phi(0)$ and $\Ht'(1)-\Ht'(0)$ vanish to working precision;
$\min_\theta\Ht'\approx0.9509$ for $(e,2)$ and $\approx0.9384$ for
$(2,e)$, with $\int_0^1\Ht'=1$ up to quadrature error; the inverse
identity $\Ht_{2,e}\circ\Ht_{e,2}=\mathrm{id}$ holds pointwise to below
$10^{-60}$; and
\[
\mu_{e,2}=-1.129041157430\ldots,\qquad
\mu_{2,e}=+1.129041157430\ldots,\qquad
\mu_{e,2}+\mu_{2,e}\approx 3\cdot10^{-14}
\]
at $4096$ phase samples (quadrature-limited), confirming
Theorem~\ref{thm:B} exactly on a branch different from the regular one.
(Incidentally, these toy mean shifts differ from the regular-branch
values $\mp1.128403\ldots$ of \cite[Table~1]{J} only in the third
decimal.)
\end{remark}


\section*{Authorship and AI Contribution Statement}
This paper is part of a five-paper series produced in a directed human--AI
collaboration. The author conceived the mixed-base phase program, built and
maintains the high-precision tetration engine used for every computation in
the series (a fork of Levenstein's \texttt{fatou.gp}; engine, data and
scripts at \url{https://github.com/Lightrunnerwastaken/gp-tetration}), and directed the research: posing the questions,
setting priorities between rounds, auditing results, and deciding which
claims met the acceptance gates. Substantial parts of the mathematical
development and of the numerical work were produced by AI systems under
that direction: Anthropic's Claude in a theory and review role, an
autonomous research agent operating the numerical pipeline, and OpenAI's
ChatGPT, which contributed to the original versions of Papers I and II and
to manuscript review. In line with prevailing journal policies, AI systems
are not listed as authors; the author assumes full and sole responsibility
for all statements and proofs, and readers should weigh this provenance in
their assessment. The five papers of the series are permanently archived on Zenodo under the concept DOIs \url{https://doi.org/10.5281/zenodo.21346803} (I), \url{https://doi.org/10.5281/zenodo.21361967} (II), \url{https://doi.org/10.5281/zenodo.21362206} (III), \url{https://doi.org/10.5281/zenodo.21363593} (IV), and \url{https://doi.org/10.5281/zenodo.21363809} (V).

\begin{thebibliography}{9}

\bibitem{J} J.~Justus, \emph{A Phase Law for Mixed-Base Tetration: Logarithmic Tails, Cocycle Structure, and Singular Asymptotics}, Paper~I of the mixed-base tetration series, preprint, July 2026. \url{https://github.com/Lightrunnerwastaken/gp-tetration}. DOI: \url{https://doi.org/10.5281/zenodo.21346803}.

\bibitem{Kneser} H.~Kneser,
\emph{Reelle analytische L\"osungen der Gleichung
$\varphi(\varphi(x))=e^{x}$ und verwandter Funktionalgleichungen},
J.~reine angew.~Math.\ \textbf{187} (1950), 56--67.

\bibitem{PC} W.~Paulsen and S.~Cowgill,
\emph{Solving $F(z+1)=b^{F(z)}$ in the complex plane},
Adv.~Comput.~Math.\ \textbf{43} (2017), 1261--1282.

\bibitem{Paulsen} W.~Paulsen,
\emph{Tetration for complex bases},
Adv.~Comput.~Math.\ \textbf{45} (2019), 243--267.

\bibitem{Apostol} T.~M.~Apostol,
\emph{Mathematical Analysis}, 2nd ed., Addison--Wesley, 1974.
(Theorems~7.6 and 7.7: Riemann--Stieltjes integration by parts and change
of variables.)

\bibitem{Rudin} W.~Rudin,
\emph{Principles of Mathematical Analysis}, 3rd ed., McGraw--Hill, 1976.
(Theorem~7.17: differentiation of uniform limits.)

\end{thebibliography}

\end{document}
